%% kapa: exam=1 round=29 year=2018 period=2 form=A subject=accounting subject_ko=회계학 seq=38 number=78
%% source: Q-Net 감정평가사 기출문제 (cst00309) artlSeq=5202808 file=extracted/1008/5202808/2018년 제29회 감정평가사 1차시험 기출문제/2018년 제29회 감정평가사 1차시험 기출문제/2교시 A형.pdf
%% conversion: kordoc 4.24.51 -> tools/gen1.py (auto); review=proofread-sonnet; flags=none
\begin{question}{78}
(주)감평은 세 종류의 제품 A, B, C를 독점 생산 및 판매하고 있다. 제품생산을 위해 사용되는 공통설비의 연간 사용시간은 총 40,000시간으로 제한되어 있다. 20x1년도 예상 자료는 다음과 같다. 다음 설명 중 옳은 것은?
\par{\small\centering
\begin{tabular}{|l|r|r|r|}
\hline
\multicolumn{1}{|c|}{구분} & \multicolumn{1}{c|}{제품 A} & \multicolumn{1}{c|}{제품 B} & \multicolumn{1}{c|}{제품 C} \\ \hline
단위당 판매가격 & ￦500 & ￦750 & ￦1,000 \\ \hline
단위당 변동원가 & ￦150 & ￦300 & ￦600 \\ \hline
단위당 공통설비사용시간 & 5시간 & 10시간 & 8시간 \\ \hline
연간 최대 시장수요량 & 2,000단위 & 3,000단위 & 2,000단위 \\ \hline
\end{tabular}\par}
\begin{choices}
  \item 제품단위당 공헌이익이 가장 작은 제품은 C이다.
  \item 공헌이익을 최대화하기 위해 생산할 제품 C의 설비 사용시간은 12,000시간이다.
  \item 공헌이익을 최대화하기 위해 생산할 총제품수량은 5,000단위이다.
  \item 공헌이익을 최대화하기 위해서는 제품 C, 제품 B, 제품 A의 순서로 생산한 후 판매해야 한다.
  \item 획득할 수 있는 최대공헌이익은 ￦2,130,000이다.
\end{choices}
\end{question}
\answer{5}
